% =====================================================================
%  日本語スライド（Beamer + LuaLaTeX）用の追加プリアンブル。
%  pandoc の --standalone に --include-in-header で足される。
%  差し替えるなら octavo.config.py に
%      'headers': {'beamer': 'my-beamer-header.tex'},
% =====================================================================

% ---- 日本語 -----------------------------------------------------------
% beamer は article 系クラスなので luatexja を明示的に読む。
% haranoaji は TeX Live 同梱なのでシステムフォントに依存しない。
\IfFileExists{luatexja-preset.sty}{%
  \usepackage{luatexja}
  \usepackage[haranoaji,deluxe]{luatexja-preset}
  % スライドの和文は明朝よりゴシックのほうが読みやすい
  \renewcommand{\kanjifamilydefault}{\gtdefault}%
}{\PackageWarningNoLine{octavo}{luatexja が無い。日本語は組めない。
   sudo apt install texlive-lang-japanese (macOS: MacTeX)}}
\renewcommand{\familydefault}{\sfdefault}

% ---- 体裁 -------------------------------------------------------------
\setbeamertemplate{navigation symbols}{}          % 右下の小さなボタンを消す
\setbeamertemplate{footline}[page number]         % ページ番号を出す（後で「n枚目」と参照できるように）
\setbeamertemplate{caption}{\insertcaption}       % 「図1:」を出さない
\setbeamerfont{frametitle}{size=\large}
\setbeamerfont{footnote}{size=\scriptsize}

% 今のsectionをフレームタイトルの上に小さく灰色で出す（長いスライドで迷子にならないように）。
% 原稿側は普通に "## 見出し" と書くだけでよい（\insertsection は beamer 側で自動的に埋まる）。
\addtobeamertemplate{frametitle}{%
  {\color{gray}\fontsize{10pt}{12pt}\selectfont\insertsection\par}%
}{}

% 箇条書きの行間を少し空ける（スライドは詰まると読めない）
\addtobeamertemplate{itemize/enumerate body begin}{\setlength{\itemsep}{0.5em}}{}

% ---- 図表 -------------------------------------------------------------
\usepackage{adjustbox}
\let\oldtabular\tabular
\let\endoldtabular\endtabular
\renewenvironment{tabular}{\adjustbox{max width=\linewidth,max height=0.62\textheight}%
                           \bgroup\oldtabular}
                          {\endoldtabular\egroup}

% 出典を小さく入れる。原稿で \source{国勢調査（2020）} と書く。
\newcommand{\source}[1]{\par\vspace{0.3em}{\scriptsize\raggedleft 出典: #1\par}}
